\documentclass[conference]{IEEEtran}
\IEEEoverridecommandlockouts
% The preceding line is only needed to identify funding in the first footnote. If that is unneeded, please comment it out.
\usepackage{cite}
\usepackage{amsmath,amssymb,amsfonts}
\usepackage{algorithmic}
\usepackage{graphicx}
\usepackage{textcomp}
\usepackage{xcolor}
\def\BibTeX{{\rm B\kern-.05em{\sc i\kern-.025em b}\kern-.08em
    T\kern-.1667em\lower.7ex\hbox{E}\kern-.125emX}}
\begin{document}

\title{Single-Cell Genomic Clustering Using scMINER with RAG-Assisted Biological Interpretation\\
{\footnotesize{Project Proposal for Clustering Algorithms course at Missouri S\&T}}
\thanks{Identify applicable funding agency here. If none, delete this.}
}

\author{\IEEEauthorblockN{1\textsuperscript{st} Utshab Kumar Ghosh}
\IEEEauthorblockA{\textit{Computer Science PhD} \\
\textit{Missouri S\&T}\\
Rolla, MO, USA \\
u.ghosh@mst.edu}
\and
\IEEEauthorblockN{2\textsuperscript{nd} Laraib Khan}
\IEEEauthorblockA{\textit{Mechanical Engineering PhD} \\
\textit{Missouri S\&T}\\
Rolla, MO, USA \\
lktmy@mst.edu}
\and
\IEEEauthorblockN{3\textsuperscript{rd} Rhea Eckman}
\IEEEauthorblockA{\textit{Systems Engineering MS} \\
\textit{Missouri S\&T}\\
St. Louis, MO, USA \\
rehry@mst.edu}
}

\maketitle

\begin{abstract}
Leaving this blank for now but including the section as it is useful
\end{abstract}

\begin{IEEEkeywords}
RAG, scRNAseq, MICA, SJARACNe, scMINER
\end{IEEEkeywords}

\section{Background and Motivation}
We selected this paper because our group is interested in exploring how clustering can support retrieval-augmented generation (RAG) for biomedical data, which is a growing research area. Single-cell transcriptomic data is a good fit for this direction because it is large, high-dimensional, and often difficult to interpret directly. The "scMINER" framework is especially suitable because it introduces a clustering approach based on mutual information that is designed to capture complex relationships among cells and identify meaningful cell populations. This gives us a strong clustering foundation for the course while also creating a natural next step where the important genes or markers from each cluster can be used to retrieve relevant biomedical literature through RAG. In this way, the project can remain focused on clustering while also extending the results toward a modern biomedical interpretation workflow. 

\section{Existing Work and Limitations}
scMINER provides an integrated framework for clustering single-cell RNA sequencing data, reconstructing regulatory networks, and identifying transcription factor and signaling drivers \cite{pan2025scminer}. Its Mutual Information-based Clustering Analysis (MICA) is designed to capture nonlinear relationships between cells and was shown to perform well compared with commonly used approaches such as Seurat, Scanpy, SC3s, scVI, and scDeepCluster. However, several aspects of the existing work leave room for further investigation.
One limitation concerns the evaluation of clustering under a fully unsupervised setting. In the original benchmarking experiments, the number of clusters for MICA and the comparison methods was selected using information from the ground-truth cell labels. For larger datasets, the Louvain resolution was varied until the predicted number of clusters matched the known number of cell populations, and the result with the highest ARI was selected when multiple settings produced the same number of clusters \cite{pan2025scminer}. This provides a controlled comparison between algorithms, but it does not completely represent exploratory single-cell analysis, where the true number of cell populations is generally unknown. Evaluating scMINER without using label information for cluster selection may therefore provide a more realistic understanding of its performance.

A second limitation is the evaluation of scMINER at very large scale. The original work demonstrated successful processing of datasets containing up to approximately 650,000 cells and reported competitive runtime and memory use \cite{pan2025scminer}. However, the authors also noted that calculating mutual-information-based cell similarities across genes is computationally demanding and that the absence of reliable ground-truth labels for atlas-scale datasets prevented accuracy evaluation at that scale \cite{pan2025scminer}. Modern single-cell atlases can contain millions of cells, creating an opportunity to investigate the tradeoff among clustering accuracy, runtime, memory use, and preservation of rare cell populations when MICA is applied beyond the scale evaluated in the original study. The authors also report that scMINER requires at least 500 cells for some analyses, which can limit its use for rare cellular populations \cite{pan2025scminer}.

A third open problem is the interpretation of the clusters after clustering. scMINER provides marker information, gene-network inference, and hidden-driver identification, but the framework does not evaluate a retrieval-grounded language-model approach for translating these outputs into literature-supported biological descriptions. Recent studies have shown that large language models can infer cell types from marker genes, while newer approaches use biomedical literature retrieval to improve the reliability of those annotations \cite{hou2024gpt4,crowley2025anndictionary,kazmi2025beyond,yang2026rag}. These studies primarily use marker-gene information as the biological input. This raises a different question for scMINER: whether its inferred regulatory drivers, when combined with conventional marker genes, can provide additional evidence for biological interpretation of clusters.
\section{Research Questions}

\section{Initial Implementation Plan}
Our plan has five steps. Everything will be done in Python. MICA and SJARACNe, the clustering and network tools behind scMINER, already run in Python, and we will rewrite the short R steps for preprocessing and driver activity in Python. The heavy runs will be done on a lab workstation with two 48 GB GPUs, 120 CPU cores and 754 GB of memory. 

\subsection{Step 1. Set up and reproduce the paper.}

All our code will live in a team GitHub repository. Every experiment will run inside a conda environment on the lab workstation, so each result can be repeated. Long runs will go in byobu sessions, which keep running if a connection drops. We will install MICA, the clustering part of scMINER, and SJARACNe, its network part, from the authors' GitHub code. Before changing anything, we will rerun MICA on two datasets used in the paper, Zeisel and PBMC14K. If our scores come close to the published ones, we can trust our setup.

\subsection{Step 2. Prepare three datasets.}
We plan to use single-cell RNA sequencing data for all three course regimes, so results can be compared across data sizes. In each dataset a sample is one cell and a variable is one gene. 


\subsubsection{Dataset 1 (ordinary).}
PBMC14K, 13,605 human blood cells and about 18,000 genes. The cell types were separated with antibodies before sequencing, so the labels do not depend on any clustering method.

\subsubsection{Dataset 2 (large scale).}
The Allen Brain Cell Atlas whole mouse brain data, about 2.3 million cells from a single 10x v3 chemistry with 32,285 genes. Labels come in levels, from 34 classes down to 338 subclasses.

\subsubsection{Dataset 3 (high dimensional).}
The Zeisel mouse cortex and hippocampus data, 3,005 cells and 19,972 genes, labeled with major cell classes and 47 subclasses. It has far more variables than samples.

For each dataset we will record the source, variables, preprocessing (quality control, normalization, log transform, gene filtering) and labels. We also have backups in case a download is impractical, such as the Cao 2020 human fetal atlas (about 4 million cells) for Dataset 2. 

\subsection{Step 3. Baselines and evaluation.}
We will compare MICA with k-means on principal components, Leiden graph clustering as used in Scanpy, and hierarchical clustering where the data is small enough. Since every dataset has labels, we will score the clusters with the adjusted Rand index (ARI) and normalized mutual information (NMI). We will also report silhouette score, run time and peak memory over several random seeds. In the paper, the number of clusters for each method was tuned until it matched the true number of cell types. We will report scores that way to compare with the paper, and also with the number of clusters chosen without labels, since that is how clustering is used in practice.

\subsection{Step 4. Extend MICA to large data.}
MICA measures how similar two cells are through the mutual information across all their genes. This captures nonlinear patterns but is expensive. The largest dataset in the paper had about 650,000 cells and ran on a 40-core CPU server, and the software has no GPU support. From our reading of the released code, it would not run on 2 million cells as it is. We plan two changes. First, we will cluster a smaller sample of cells chosen to keep rare cell types, then assign every other cell to its nearest cluster. Second, we will compute the mutual information in GPU batches so that larger samples, or even the full data, become practical. We will then measure accuracy, run time and memory trade off, and check whether small cell groups get lost.

\subsection{Step 5. Interpret clusters with retrieval (RAG)}
For each cluster we will collect its top marker genes and, where the network step runs, the hidden driver genes reported by scMINER. These genes will serve as search queries over marker gene databases such as CellMarker 2.0 and PanglaoDB, and over PubMed abstracts. A language model running locally on our GPUs will read the retrieved text and suggest a cell type name for the cluster, citing its sources. We will compare these names with the known labels. This checks whether the clusters make biological sense and gives a short, readable description of each one. 

\section{Related Work and Expected Contribution}
Recent research in single-cell transcriptomics has focused on both computational clustering and automated biological interpretation. Pan et al. introduced scMINER in 2025 as a mutual-information-based framework that combines cell clustering with transcription-factor and signaling-network inference \cite{pan2025scminer}. Its clustering component, MICA, uses mutual information to capture nonlinear relationships between cells and was evaluated against several commonly used single-cell clustering methods. In addition to identifying cell populations, scMINER provides marker genes and regulatory-driver information that can support downstream biological analysis.

Large language models have also been increasingly studied for the interpretation of single-cell data. Hou and Ji evaluated GPT-4 for cell-type annotation using marker-gene lists and showed that language models can assist in assigning biological identities to cell clusters \cite{hou2024gpt4}. More recently, Crowley et al. introduced AnnDictionary and benchmarked several large language models for de novo cell-type annotation using marker genes derived from unsupervised clustering \cite{crowley2025anndictionary}. Their results showed that LLM-based annotation can be useful, although performance varies across models, datasets, and cell populations.

More recent studies have incorporated retrieval-based methods to provide external biological evidence during cell-type annotation. Kazmi et al. explored a retrieval-augmented generation approach in which marker-gene information was supplemented with relevant biomedical literature before generating cell-type annotations \cite{kazmi2025beyond}. Yang et al. similarly developed a RAG-based framework that combined lexical and semantic retrieval, query expansion, reranking, and ontology-based evaluation for single-cell annotation \cite{yang2026rag}. These studies indicate that incorporating retrieved biomedical evidence can improve the grounding and reliability of language-model-generated annotations.

\section{Rough Timeline}
\begin{itemize}
\item{Late September.}
Set up the software, download and preprocess the data, and reproduce the paper on Zeisel and PBMC14K. 
\item{October.}
Run the baselines on all three datasets and start the large-scale runs.
\item{November.}
Build and test the large-data extension and the retrieval step.
\item{December.}
Final experiments, report and presentation. 
\end{itemize}

\section*{Acknowledgment}

The preferred spelling of the word ``acknowledgment'' in America is without 
an ``e'' after the ``g''. Avoid the stilted expression ``one of us (R. B. 
G.) thanks $\ldots$''. Instead, try ``R. B. G. thanks$\ldots$''. Put sponsor 
acknowledgments in the unnumbered footnote on the first page.

\bibliographystyle{IEEEtran}
\bibliography{references}
\vspace{12pt}

\end{document}
